% !TEX root = ../main.tex
\chapter*{DANH MỤC KÝ HIỆU}
\addcontentsline{toc}{chapter}{DANH MỤC KÝ HIỆU}

\begin{tabular}{@{}L{3.5cm}L{9.8cm}@{}}
    \textbf{Ký hiệu} & \textbf{Ý nghĩa và đơn vị} \\
    $E$, $B$ & Hệ tọa độ quán tính cục bộ và hệ tọa độ thân \\
    $R$ & Ma trận chuyển tọa độ từ hệ $B$ sang hệ $E$ \\
    $\phi,\theta,\psi$ & Góc roll, pitch, yaw (rad) \\
    $p,q,r$ & Vận tốc góc theo ba trục thân (rad/s) \\
    $\Omega_i$ & Tốc độ góc rôto thứ $i$ (rad/s) \\
    $f_i$, $T$ & Lực đẩy rôto thứ $i$ và tổng lực đẩy (N) \\
    $Q_i$ & Độ lớn mômen cản khí động rôto thứ $i$ (N$\cdot$m) \\
    $c_T$, $c_Q$ & Hệ số lực đẩy và hệ số mômen cản theo bình phương tốc độ rôto \\
    $l$, $a$ & Chiều dài tay đòn và hình chiếu $a=l/\sqrt{2}$ (m) \\
    $J_x,J_y,J_z$ & Mômen quán tính quanh ba trục thân (kg$\cdot$m$^2$) \\
    $\tau_x,\tau_y,\tau_z$ & Mômen điều khiển quanh ba trục thân (N$\cdot$m) \\
    $d_x,d_y,d_z$ & Các thành phần mômen nhiễu tổng hợp (N$\cdot$m) \\
    $m$, $g$ & Khối lượng (kg) và gia tốc trọng trường (m/s$^2$) \\
    $T_s$ & Chu kỳ lấy mẫu hoặc cập nhật thuật toán (s) \\
    $K_\phi,K_\theta$ & Hệ số tỉ lệ của vòng góc (s$^{-1}$) \\
    $K_{P,j},K_{I,j},K_{D,j}$ & Các hệ số PID vận tốc góc của kênh $j$ \\
    $\alpha_f$, $\tau_f$ & Hệ số phối hợp và hằng số thời gian bộ lọc bổ sung (s đối với $\tau_f$) \\
    $\kappa$ & Tỉ số $c_Q/c_T$ (m) \\
    $\varepsilon_\chi$ & Dải sai lệch cho phép của góc $\chi$ (rad) \\
    $M_{p,\chi}$, $t_{s,\chi}$ & Độ quá điều chỉnh (\%) và thời gian xác lập (s) \\
\end{tabular}
